\documentclass{article}
\usepackage[utf8]{inputenc}
\usepackage[T1]{fontenc}
\usepackage[french]{babel}
\usepackage{amsmath,amssymb}
\usepackage{geometry}
\geometry{margin=2cm}

\title{Compte rendu du projet de CoCa}
\author{Pourageaud Matthieu / Vasseur Cengiz}
\date{\today}

\begin{document}

\maketitle

\section{Introduction}

On définit dans notre structure les variables booléennes :

\medskip

\noindent
$x_{n,p,h}$ : le sommet $n$ à la position $p$ a une hauteur de pile $h$.\\
$y_{p,h,a}$ : à la position $p$, il y a un paquet de valeur $a$ à l'étage $h$ de la pile.

\medskip

On définit aussi :
\begin{itemize}
  \item $s$ : le sommet de départ,
  \item $e$ : le sommet d'arrivée,
  \item $length$ : la longueur du chemin,
  \item $Op$ : l'ensemble d'opérations
  \[
    \{\,4 \to 4,\; 6 \to 6,\; 4 \uparrow 44,\; 4 \uparrow 46,\; 6 \uparrow 64,\; 6 \uparrow 66,\; 44 \downarrow 4,\; 46 \downarrow 4,\; 64 \downarrow 6,\; 66 \downarrow 6\,\},
  \]
  \item $h_{\max}$ : la hauteur maximale de la pile,
\end{itemize}

Soit le réseau \[R = ⟨G = (V, E), PUSH_{a}^{b}, POP_{a}^{b}, T_{a} | a,b \in {4, 6}⟩\]\\
Problème : Existe-t-il un chemin simple et valide de $s$ à $e$ de longueur $length$ dans le réseau $R$ ?

\subsection{Début et fin de chemin}

Début (s, V, hmax) :
\[
x_{s,0,0} \land y_{0,0,4} \land \lnot y_{0,0,6}
\land \bigwedge_{n\in V} \lnot x_{n,0,0}
\land \bigwedge_{h=1}^{h_{\max}} \bigwedge_{n\in V} \lnot x_{n,0,h}.
\]

Fin (e, length, V, hmax) :
\[
x_{e,\;length,\;0} \land y_{length,0,4} \land \lnot y_{length,0,6}
\land \bigwedge_{n\in V} \lnot x_{n,\,length,\,0}
\land \bigwedge_{h=1}^{h_{\max}} \bigwedge_{n\in V} \lnot x_{n,\,length,\,h}.
\]

\subsection{Chemin simple et valide}
\[
\bigwedge_{p=0}^{l} \bigwedge_{h=0}^{h_{\max}} \bigwedge_{n\in V} \Bigl( \text{Unique}\bigl(x_{n,p,h}\bigr).
\land
\bigl( x_{n,p,h} \Rightarrow
\bigvee_{\substack{n' \in V \\ n' \neq n}} \bigvee_{op \in Op}
\bigl(\text{Transmission}\ \lor\ \text{Encapsule}\ \lor\ \text{Decapsule}\bigr) \bigr) \Bigr).
\]

\subsubsection*{Transmission (op = $4\to 4$ ou $6\to 6$)}
\[
\text{Transmission (p, h, a, n\')}:\qquad
y_{p,h,a} \land x_{n',\,p+1,\,h} \land y_{p+1,\,h,\,a}
\iff y_{p,h,a} \land y_{p+1,h,a}.
\]

\subsubsection*{Encapsule (op = $4\uparrow44$, $4\uparrow46$, $6\uparrow64$, $6\uparrow66$)}
\[
\text{Encapsule (n, p, h, n\') si $h < h_{\max}$}:\qquad
x_{n,p,h} \Rightarrow
\bigl( y_{p,h,a} \land x_{n',\,p+1,\,h+1} \land y_{p+1,\,h+1,\,b} \land y_{p+1,\,h,\,4} \bigr).
\]

\subsubsection*{Decapsule (op = $44\downarrow4$, $46\downarrow4$, $64\downarrow6$, $66\downarrow6$)}
\[
\text{Decapsule (n, p, h, n\') si $h>0$}:\qquad
x_{n,p,h} \Rightarrow
\bigl( y_{p,h,a} \land x_{n',\,p+1,\,h-1} \land y_{p+1,\,h-1,\,b} \land y_{p+1,\,h,\,4} \bigr).
\]

\subsection{Pile cohérente}

\paragraph{Partie 1 (length, V, hmax) :}
Pour chaque position $p$ :
\[
\bigwedge_{p=0}^{length}\Biggl(
\bigl(\bigvee_{n\in V} x_{n,p,0}\bigr) \Rightarrow \lnot\bigl(y_{p,0,4}\land y_{p,0,6}\bigr) \land
\bigwedge_{h=1}^{h_{\max}} \Bigl(
\bigl(\bigvee_{n\in V} x_{n,p,h}\bigr) \Rightarrow 
\bigl((y_{p,h,4}\land \lnot y_{p,h,6}) \lor (\lnot y_{p,h,4}\land y_{p,h,6})\bigr)
\Bigr). \Biggr) 
\]

\paragraph{Partie 2 (length, V, hmax) :}
\[
\bigwedge_{n\in V}\bigwedge_{p=0}^{length}\bigwedge_{h=0}^{h_{\max}}
\bigl( x_{n,p,h} \Rightarrow (H_{\text{grand}} \land H_{\text{petit}}) \bigr).
\]

\noindent
Soit $H_{\text{grand}}$ :
\[
H_{\text{grand}}:\qquad
\bigwedge_{h'=h+1}^{h_{\max}} \Bigl( \lnot\bigl(x_{n,p,h} \land y_{p,h',4}\bigr) \land \lnot\bigl(x_{n,p,h} \land y_{p,h',6}\bigr) \Bigr).
\]

\noindent
Soit $H_{\text{petit}}$ :
\[
H_{\text{petit}}:\qquad
\bigwedge_{h'=0}^{h} \Bigl( \lnot x_{n,p,h} \lor y_{p,h',4} \lor y_{p,h',6} \Bigr).
\]

\subsection{Unique}

\paragraph{Au plus:}
Soit Au plus (length, V, hmax) :
\[
\bigwedge_{n\in V}\bigwedge_{p=0}^{length}\bigwedge_{h=0}^{h_{\max}\bigwedge_{n' in V n \neq n'}}
\bigl( \lnot x_{n,p,h} \land \lnot x_{n',p,h} \bigr).
\]

\paragraph{Au moins :}
Soit Au moins (length, V, hmax) :
\[
\bigwedge_{n\in V}\bigwedge_{p=0}^{length}\bigwedge_{h=0}^{h_{\max}}.
x_{n,p,h}.
\]

\paragraph{Unique :}
Soit Unique (length, V, hmax) :
\[
\bigwedge_{n\in V}\bigwedge_{p=0}^{length}\bigwedge_{h=0}^{h_{\max}}.
\text{Au moins (n, p, h)} \land \text{Au plus (n, p, h)}.
\]

\section{Formules}

La formule finale sera alors :

\[
\text{Début} \land \text{Fin} \land \text{Chemin simple} \land \text{Chemin valide} \land \text{Pile cohérente}
\]


\end{document}
